%%%%%%%%%%%%%%%%%%%%%%%%%%%%%%%%%
%
%       EXERCÍCIO
%
%%%%%%%%%%%%%%%%%%%%%%%%%%%%%%%%%

\ifdefstring{\atividade}{prova}{%
    \renewcommand{\valorquestao}{\ValorQ\ pontos}
}%

\begin{exercicioBanco}[\valorquestao]
    Para se avaliar o nível de tensão ocasionada por exames escolares, doze alunos foram escolhidos e sua pulsação medida antes e depois do exame.

    \begin{center}
    \begin{tabular}{|c|cccccccccccc|}
        \hline
        \multirow{2}{*}{Instante da medição} & \multicolumn{12}{c|}{Estudante} \\
        \cline{2-13}
        & 1 & 2 & 3 & 4 & 5 & 6 & 7 & 8 & 9 & 10 & 11 & 12 \\
        \hline
        Antes  & 87 & 78 & 85 & 93 & 76 & 80 & 82 & 77 & 91 & 74 & 76 & 79 \\
        Depois & 83 & 84 & 79 & 88 & 75 & 81 & 74 & 71 & 78 & 73 & 76 & 71 \\
        \hline
    \end{tabular}
    \end{center}

    Faça um teste, com nível de significância de \(1\%\), para verificar se existe maior tensão (isto é, maior pulsação) antes da realização dos exames. Indique as suposições necessárias.
\end{exercicioBanco}

